\section{Introduction}
\label{sec:intro}

Recent work has demonstrated that large language models (LLMs) can serve
as algorithm designers---systems that write code to solve optimization
problems, then iterate on that code based on evaluation
feedback~\cite{fawzi2024funsearch,deepmind2025alphaevolve,ye2024reevo}.
These results are striking, but they share a common evaluation pattern:
generated code is tested on instances drawn from the same distribution
used during search. For engineering optimization, where the objective
landscape changes with physical parameters, this conflates
hyperparameter tuning with algorithmic discovery.

We present FunWake, an evaluation framework for LLM agents that write
wind farm layout optimizers. The key design features are:
\begin{enumerate}[nosep]
    \item \textbf{Held-out evaluation with different physics.} A
    held-out farm differs from training in turbine model, polygon
    geometry, turbine count, and wind resource. The LLM never sees
    held-out AEP---only a binary feasibility signal.
    \item \textbf{Controlled autonomy levels.} In \emph{full-optimizer}
    mode, the LLM writes arbitrary \texttt{optimize()} functions. In
    \emph{schedule-only} mode, it writes only a 4-parameter schedule
    function controlling learning rate, penalty weight, and Adam
    momentum within a fixed optimization skeleton.
    \item \textbf{Multi-model, multi-scaffold comparison.} We evaluate
    Claude Code (frontier), Qwen~2.5 Coder~32B, and Llama~3.3~70B
    (open-source) through both a rich CLI agent and a minimal
    structured-output interface.
\end{enumerate}

Three findings emerge from this framework:

\paragraph{Scaffolding dominates model capability in constrained
settings.} Open-source models that never read the optimization
codebase---because their minimal agent interface discourages
exploration---produce schedules stuck 13\,GWh below baseline.
Injecting the same codebase context that Claude Code reads via its
tools produces a +20\,GWh improvement, closing the gap entirely.
The performance difference previously attributed to model capability
was largely information asymmetry.

\paragraph{The optimal autonomy level is model-dependent.}
Claude Code's schedule-only mode achieves the best held-out
generalization (+24.8\,GWh over baseline), outperforming its
full-optimizer mode (+21.8\,GWh). For open-source models, the pattern
reverses: full-optimizer mode generalizes better (+5\,GWh) while
schedule-only fails to produce feasible held-out layouts. Constraining
the action space helps frontier models but hinders smaller ones.

\paragraph{Training performance misleads without held-out evaluation.}
The script with the highest training AEP (5549.5\,GWh, +8.8 over
baseline) has a held-out score \emph{below} baseline. The script with
the best held-out generalization has a lower training score. Any
evaluation that reports only training AEP would draw the wrong
conclusions about which optimizer is best.

Our contributions are: (1)~an open-source benchmark for LLM-generated
optimization code with held-out engineering evaluation;
(2)~a controlled experiment showing that agent scaffolding explains
more variance than model scale in constrained optimization settings;
and (3)~evidence that the optimal LLM autonomy level depends on model
capability, necessitating benchmark designs that test multiple levels.
